% !TeX root = main.tex
\clearpage
\subsection{A5: Plan de migración y matriz de riesgos}
\label{sec:fase5-a5}

\subsubsection{Alcance y reglas de transición}

El plan desarrolla las secciones 4.5 y 4.6 de la consigna: orden de extracción
justificado, criterio de hecho por etapa, reversión planificada en menos de
un minuto y matriz de probabilidad por impacto con mitigaciones. Se apoya
en ADR-006--010, contratos A3 y modelo A4. Las fuentes editables se conservan
en \path{cna/docs/a5/}.

BookStack conserva su código, contenido, permisos y MariaDB. La integración
se limita a REST y Webhooks; no se entregan credenciales de MariaDB a los
servicios nuevos. Se habilita una funcionalidad por vez, después de probarla
en paralelo, y se mantiene disponible su alternativa nativa. Cerrar una
entrada no implica borrar colas, eventos, trabajos ni archivos.

\textbf{Estado: plan de diseño.} Los criterios de hecho siguientes son
condiciones de salida para ejecutar en Fases 6 y 7, no pruebas aprobadas.
La duración de reversión y la recuperación de datos no se han medido.
El cierre documental de Fase 5 no equivale a cumplir los criterios globales
CA-1--CA-7 del trabajo completo.

\subsubsection{Orden de extracción y criterio de hecho}

El orden relativo es \textbf{S2 auditoría, S1 búsqueda, S3 exportación y S4
medios}. Primero se agregan registros y proyecciones sin reemplazar el
contenido; después se incorporan trabajos y objetos con referencias
duraderas. La prioridad se refiere al impacto funcional y al costo de
reversión en esta instalación. No implica que se tolere perder auditoría.

\begin{table}[H]
\centering\small
\begin{tabularx}{\textwidth}{>{\raggedright\arraybackslash}p{1.1cm} >{\raggedright\arraybackslash}p{2.7cm} >{\raggedright\arraybackslash}X}
\toprule
\textbf{Paso} & \textbf{Alcance} & \textbf{Justificación y condición de habilitación} \\
\midrule
M0 & Preparación & Proteger línea base, validar recursos, backup/restauración, red, secretos y variantes de gateway antes de abrir rutas. \\
M1 & S2 auditoría & Registro paralelo sin intervenir la edición; habilitar después de deduplicar, recuperar pendientes y probar filtros. \\
M2 & S1 búsqueda & Proyección reconstruible; abrir tras verificar autorización, reconciliación y actualización menor que 5 s. \\
M3 & S3 exportación & Jobs y objetos con efectos en curso; abrir POST después de probar reserva recuperable, estado y descarga. \\
M4 & S4 medios & Objetos potencialmente referenciados por contenido; exige retención y restauración antes de nuevas referencias. \\
M5 & Consolidación & Cotejar despliegue, criterios globales, pruebas y reversión medida antes de la entrega final. \\
\bottomrule
\end{tabularx}
\caption{Orden de migración de A5, de menor a mayor riesgo funcional, con preparación y consolidación.}
\label{tab:a5-orden}
\end{table}

\clearpage
\paragraph{M0: preparar sin alterar la línea base.}
Guardar backup consistente de base, archivos, configuración y secretos
necesarios para restaurar, fuera de Git y con acceso restringido. Probar la
restauración en un entorno aislado. Verificar RAM del anfitrión, recursos
de ambas VM, red privada, certificados y puertos; preparar límites,
variantes de gateway y probes nativos antes de habilitar VM-2.
\textbf{Hecho:} restauración evidenciada; login, creación de libro y edición
intactos; red y restricciones comprobadas; presupuesto medido y entrada del
puente persistida antes de 202, recuperable ante caída del broker.
Responsable previsto: infraestructura, con revisión de otro integrante.

\paragraph{M1: auditoría con S2.}
Crear base y rol propios y consumidor durable independiente. Contrastar
eventos aceptados con actividad nativa sin cambiar páginas.
\textbf{Hecho:} filtros de A3 correctos; un efecto por \texttt{event\_id};
S2 detenido dos minutos durante veinte ediciones recupera los veinte eventos
publicados; un mensaje inválido conserva su causa en DLQ; el monolito sigue
operativo. La prueba diferencia eventos aceptados de Webhooks que no
llegaron al puente. Responsable: integración/eventos; CA-5/CA-6.

\paragraph{M2: búsqueda con S1.}
Crear índice y rol propios, establecer consumidor/cursor antes de cargar
páginas y libros por API, conservar cambios concurrentes y reconciliar
por entidad y revisión antes de publicar la búsqueda.
\textbf{Hecho:} diez ediciones controladas aparecen cada una en menos de
cinco segundos; creación, actualización y borrado son coherentes; la carga
inicial con ediciones no pierde cambios; no hay contenido fuera del alcance
autorizado; sin coincidencias se devuelve 200 vacío; S1 caído produce 503
JSON mientras el monolito permanece disponible. Las diez ediciones son un
ensayo propuesto, no otro umbral docente. Responsable: búsqueda.
Referencias: CA-3/CA-5 y F6-E4/P36; rendimiento pendiente de Fase 7.

\paragraph{M3: exportación con S3.}
Preparar Redis persistente, worker, outbox y bucket; probar página/libro
en PDF/Markdown antes de admitir solicitudes nuevas.
\textbf{Hecho:} 202 con job consultable y archivo válido; idempotencia y 409
con cuerpo distinto; Redis inaccesible antes de aceptar devuelve 503;
worker interrumpido permite recuperar el mismo job sin bloqueo indefinido;
fallo permanente queda registrado; URL firmada abre por el gateway y el
outbox conserva publicaciones pendientes. Validar AOF y recuperación del
volumen antes de atribuir durabilidad. Responsable: exportación/persistencia.
Referencias: F6-E6/P38, F6-E7/P39 y F7-E3/P51.

\paragraph{M4 y M5: medios y consolidación.}
S4 se habilita después de probar original, miniatura proporcional, límites,
metadatos, autorización, publicación recuperable y restauración aislada.
Los objetos deben seguir accesibles tras reiniciar con el volumen; las
cargas nativas se conservan. Responsable: medios/almacenamiento;
referencia F6-E8/P40. M5 exige demostrar CA-1--CA-7, ejecutar las pruebas
de Fases 6/7, cotejar diagramas y despliegue y medir la reversión antes de
etiquetar la entrega. Estos pasos permanecen pendientes de ejecución.

\clearpage
\subsubsection{Reversión por ruta: objetivo menor que un minuto}

Volver al monolito significa cerrar la entrada nueva y continuar por su
función nativa, que nunca se elimina. Las respuestas y autenticación de
\texttt{/v2/*} no equivalen a las de BookStack: no se redirige un POST a una
API incompatible ni se presenta HTML como JSON de A3. El consumidor recibe
503 JSON uniforme, correlación y un enlace alternativo; el usuario continúa
con su sesión nativa autorizada.

\begin{table}[H]
\centering\small
\begin{tabularx}{\textwidth}{>{\raggedright\arraybackslash}p{2.3cm} X X}
\toprule
\textbf{Selector} & \textbf{Acción y alternativa nativa} & \textbf{Conservación necesaria} \\
\midrule
\texttt{audit} & Retirar consulta afectada; abrir actividad nativa con sus permisos. & Conservar consumidor/base CNA: la actividad nativa no contiene necesariamente eventos exclusivos de servicios. \\
\texttt{search} & Retirar GET nuevo; usar búsqueda nativa con término codificado. & Mantener consumidor e índice para corregir y reconciliar. \\
\texttt{exports} & Bloquear nuevos POST; exportar desde página/libro nativos. & Mantener GET de jobs, workers recuperables y descargas de trabajos ya aceptados. \\
\texttt{media} & Bloquear nuevos POST; usar editor y cargas nativas. & Mantener GET de metadatos, imágenes y miniaturas ya referenciadas. \\
\bottomrule
\end{tabularx}
\caption{Reversión funcional por selector, sin eliminar datos ni simular compatibilidad entre APIs.}
\label{tab:a5-alternativas}
\end{table}

Antes de habilitar rutas se prepararán dos variantes de Nginx por selector,
\texttt{cna} y \texttt{native}, con un enlace activo por funcionalidad,
backup conocido válido y bloqueo de cambios concurrentes. La variante
nativa restringirá nuevas admisiones según la tabla y conservará el resto
del gateway. Los archivos se implementarán en Fase 6: no existen como
configuración instalada durante A5.

Cada variante debe validarse en la configuración completa. El operador
prepara acceso y probes autorizados antes de iniciar el reloj; conserva el
destino previo, cambia el enlace de forma atómica, ejecuta
\texttt{nginx -t} y solo si es correcto realiza
\texttt{systemctl reload nginx}. Si la prueba falla, restaura el enlace
previo sin recargar. Si falla la recarga o el comportamiento no cambia,
restaura, valida y recarga la variante conocida válida, registrando fallo.
No se detiene ni reinicia Nginx como mecanismo normal de reversión.

\clearpage
\begin{table}[H]
\centering\small
\begin{tabularx}{\textwidth}{>{\raggedright\arraybackslash}p{2.5cm} X}
\toprule
\textbf{Presupuesto} & \textbf{Acción prevista} \\
\midrule
0--10 s & Capturar estado anterior y cambiar atómicamente el selector preparado. \\
10--20 s & Validar configuración efectiva; si falla, restaurar y no recargar. \\
20--30 s & Recargar sin detener servicio y comprobar cambio con conexiones nuevas. \\
30--55 s & Verificar respuesta de la ruta retirada, operación nativa autorizada y lecturas previas que deban conservarse. Registrar fin. \\
55--59 s & Margen operativo; completar en menos de 60 s. \\
\bottomrule
\end{tabularx}
\caption{Presupuesto de tiempo de reversión. Es un objetivo de diseño, no una duración medida.}
\label{tab:a5-tiempo}
\end{table}

La medición comienza al ordenar la reversión, con sesión y configuración
preparadas, y termina al comprobar la ruta y la función nativa. El tiempo
de detección previo se informa aparte. Enviar la señal de recarga no prueba
que la nueva configuración atienda solicitudes: deben usarse conexiones
nuevas. Login 200 tampoco demuestra una búsqueda, exportación o carga
funcional; se realizará la acción equivalente con una sesión autorizada.

\textbf{Aceptación:} registrar \(t_1-t_0<60\,s\) por selector, respuestas
sanitizadas y operación nativa. Repetir el ensayo para cada funcionalidad.
Si se excede el umbral, registrar incumplimiento y ajustar el procedimiento.
No incluir instalaciones, backups completos ni creación de VM en esa ventana.
El runbook detallado, incluido un ejemplo de comandos condicionado a tener
los snippets preparados, está en \path{cna/docs/a5/reversion.md}.

Cerrar POST no recupera un almacén caído ni traslada archivos a BookStack.
Los GET previos solo podrán responder si sus dependencias están disponibles.
Los jobs aceptados se drenan o recuperan; los buckets y eventos no se borran.
La recuperación completa del índice, trabajos u objetos puede tardar más y
se mide aparte. Si VM-1 o el gateway completos están caídos, este rollback
no puede ejecutarse y no se le atribuye el objetivo de un minuto.

\subsubsection{Evaluación de riesgos}

La matriz comienza en la Tabla~\ref{tab:a5-riesgos}, dividida en partes,
con cinco riesgos exigidos y ocho adicionales. Probabilidad e impacto se
estiman de 1 a 3; P/I/R muestra ambos valores y su producto: 1--2 bajo,
3--4 medio y 6--9 alto. El residual supone medidas aplicadas y verificadas;
no es el riesgo actual ni una probabilidad estadística. Sus factores y estado
se conservan en el CSV.

R07 sigue alto sin recuperación validada del origen. R05 conserva impacto
alto porque no se agrega alta disponibilidad al gateway.

% Generado por generar-matriz.py. Editar el CSV y regenerar.
\clearpage
\begin{table}[H]
\centering\small
\setlength{\tabcolsep}{4pt}
\renewcommand{\arraystretch}{1.10}
\caption{Matriz de riesgos de A5, parte 1/4: probabilidad, impacto, mitigaciones y puntaje residual estimado.}\label{tab:a5-riesgos}
\begin{tabularx}{\textwidth}{>{\raggedright\arraybackslash}p{0.75cm} >{\raggedright\arraybackslash}p{2.8cm} >{\raggedright\arraybackslash}p{1.8cm} >{\raggedright\arraybackslash}X}
\toprule
\textbf{ID} & \textbf{Riesgo} & \textbf{P/I/R} & \textbf{Mitigación, señal y verificación} \\
\midrule
R01 & Consistencia eventual visible & 3/2/6\par Residual: 2 & Consumidor durable; cursor antes de carga inicial; reconciliar por API y serializar por entidad; informar retraso y retirar búsqueda si incumple.\par\textbf{Señal:} Edición no visible antes de 5 s o proyección divergente.\par\textbf{Rol:} Búsqueda. \textbf{Verificar:} Medir diez ediciones individualmente; probar creación y borrado; repetir tras reentrega. \\[4pt]
R02 & Presión de RAM en anfitrión y VM & 3/3/9\par Residual: 4 & Medir recursos reales; liberar memoria antes de VM-2; presupuestar límites y reservar margen para SO y monolito; habilitar dependencias por etapas.\par\textbf{Señal:} OOM o presión sostenida; el monolito degrada al agregar servicios.\par\textbf{Rol:} Infraestructura. \textbf{Verificar:} Capturar memoria y docker stats en reposo/carga; detener admisiones si se excede el presupuesto. \\[4pt]
R03 & Latencia de red entre VM y API & 2/2/4\par Residual: 2 & Red privada; timeout; backoff con jitter; circuito; límites de concurrencia; no saturar PHP-FPM.\par\textbf{Señal:} Timeouts o colas crecientes; latencia inducida de 2 s.\par\textbf{Rol:} Integración. \textbf{Verificar:} Ensayo X5; medir p95 y recuperación sin saturación; contrastar con línea base. \\[4pt]
R04 & Secretos compartidos o filtrados & 2/3/6\par Residual: 3 & Roles y credenciales por servicio; secretos externos a Git; permisos mínimos; redacción de tokens y cookies; rotación y revocación.\par\textbf{Señal:} Token en un log o acceso cruzado entre bases/servicios.\par\textbf{Rol:} Seguridad. \textbf{Verificar:} Probar rechazo con rol ajeno y token revocado; revisar evidencias sanitizadas y rotación. \\[4pt]
\bottomrule
\end{tabularx}
\end{table}
\clearpage
\begin{table}[H]
\centering\small
\setlength{\tabcolsep}{4pt}
\renewcommand{\arraystretch}{1.10}
\caption{Matriz de riesgos de A5, parte 2/4: probabilidad, impacto, mitigaciones y puntaje residual estimado.}\label{tab:a5-riesgos-2}
\begin{tabularx}{\textwidth}{>{\raggedright\arraybackslash}p{0.75cm} >{\raggedright\arraybackslash}p{2.8cm} >{\raggedright\arraybackslash}p{1.8cm} >{\raggedright\arraybackslash}X}
\toprule
\textbf{ID} & \textbf{Riesgo} & \textbf{P/I/R} & \textbf{Mitigación, señal y verificación} \\
\midrule
R05 & Punto único de falla en gateway y VM-1 & 2/3/6\par Residual: 3 & Mantener configuración válida respaldada; nginx -t y reload; monitorear; preparar recuperación del nodo. Aceptar falta de HA en laboratorio.\par\textbf{Señal:} Gateway o VM-1 no responde; falla de recarga.\par\textbf{Rol:} Infraestructura. \textbf{Verificar:} Probar recuperación de configuración y medir MTTR; no prometer rollback de 60 s con nodo caído. \\[4pt]
R06 & Pérdida por retención o fallo del broker & 2/3/6\par Residual: 3 & JetStream en volumen; límites y retención suficientes; inbox/outbox persistente; ACK tras commit; transferencia DLQ confirmada y respaldos.\par\textbf{Señal:} Stream cerca de cuota o evento aceptado no publicado.\par\textbf{Rol:} Eventos. \textbf{Verificar:} Detener S2 durante veinte ediciones; verificar pendientes y recuperación; probar broker caído con outbox. \\[4pt]
R07 & Webhook que nunca llega al puente & 2/3/6\par Residual: 6 & Monitorear errores de origen; validar recuperación de origen sin cambiar BookStack; reconciliar índice por API. No afirmar auditoría histórica completa sin esa validación.\par\textbf{Señal:} Edición nativa sin entrada persistida en el puente.\par\textbf{Rol:} Integración. \textbf{Verificar:} Interrumpir tramo de origen y contrastar acciones; validar qué datos se recuperan. Riesgo sigue alto mientras falte procedimiento probado. \\[4pt]
R08 & Exposición de contenido sin autorización & 2/3/6\par Residual: 3 & Limitar alcance compartido de laboratorio; validar credenciales antes de devolver texto o archivos; roles y red privada; no exponer puertos de servicio.\par\textbf{Señal:} Cliente obtiene página o archivo fuera de su alcance.\par\textbf{Rol:} Seguridad. \textbf{Verificar:} Casos negativos de búsqueda/exportación/medios; comprobar acceso externo únicamente por gateway. \\[4pt]
\bottomrule
\end{tabularx}
\end{table}
\clearpage
\begin{table}[H]
\centering\small
\setlength{\tabcolsep}{4pt}
\renewcommand{\arraystretch}{1.10}
\caption{Matriz de riesgos de A5, parte 3/4: probabilidad, impacto, mitigaciones y puntaje residual estimado.}\label{tab:a5-riesgos-3}
\begin{tabularx}{\textwidth}{>{\raggedright\arraybackslash}p{0.75cm} >{\raggedright\arraybackslash}p{2.8cm} >{\raggedright\arraybackslash}p{1.8cm} >{\raggedright\arraybackslash}X}
\toprule
\textbf{ID} & \textbf{Riesgo} & \textbf{P/I/R} & \textbf{Mitigación, señal y verificación} \\
\midrule
R09 & Jobs duplicados o bloqueados en running & 2/2/4\par Residual: 2 & Idempotency-Key; creación atómica; reserva renovable y recuperable; clave estable de objeto; intentos acotados y failed con causa.\par\textbf{Señal:} Reserva vencida sin recuperación o más de un job por misma solicitud.\par\textbf{Rol:} Exportación. \textbf{Verificar:} Matar worker; recuperar mismo job; repetir clave; inspeccionar estados y objeto final. \\[4pt]
R10 & Objetos perdidos o referencias rotas al revertir & 2/3/6\par Residual: 3 & Conservar GET y buckets al cerrar POST; copias y restauración aislada; no reemplazar URLs nativas masivamente; firma S3 accesible por gateway.\par\textbf{Señal:} Imagen referenciada o descarga de job deja de abrir.\par\textbf{Rol:} Almacenamiento. \textbf{Verificar:} Abrir objetos previos después de rollback/reinicio; probar restauración y vencimiento de URL. \\[4pt]
R11 & Cambio de IP o red del laboratorio & 3/2/6\par Residual: 2 & Usar red Solo host entre nodos y nombres configurables; registrar direcciones vigentes; reservar o fijar IP tras revisar DHCP; actualizar hosts sin cambiar contratos.\par\textbf{Señal:} Cambiar Wi-Fi rompe conexión o resolución.\par\textbf{Rol:} Infraestructura. \textbf{Verificar:} Repetir conectividad por ens37 y resolución de dominios al cambiar la red física. \\[4pt]
R12 & Reversión lenta o incompatible & 2/3/6\par Residual: 3 & Preparar variantes por ruta; cambio atómico con bloqueo; validar y recargar; probes nuevos y función nativa; conservar jobs y objetos. No redirigir POST a API incompatible.\par\textbf{Señal:} No termina antes de 60 s o la API retirada sigue recibiendo nuevas solicitudes.\par\textbf{Rol:} Operación. \textbf{Verificar:} Cronometrar cada selector; probar fallo de validación y restauración; guardar respuestas y operación nativa. \\[4pt]
\bottomrule
\end{tabularx}
\end{table}
\clearpage
\begin{table}[H]
\centering\small
\setlength{\tabcolsep}{4pt}
\renewcommand{\arraystretch}{1.10}
\caption{Matriz de riesgos de A5, parte 4/4: probabilidad, impacto, mitigaciones y puntaje residual estimado.}\label{tab:a5-riesgos-4}
\begin{tabularx}{\textwidth}{>{\raggedright\arraybackslash}p{0.75cm} >{\raggedright\arraybackslash}p{2.8cm} >{\raggedright\arraybackslash}p{1.8cm} >{\raggedright\arraybackslash}X}
\toprule
\textbf{ID} & \textbf{Riesgo} & \textbf{P/I/R} & \textbf{Mitigación, señal y verificación} \\
\midrule
R13 & Disco lleno o pérdida de volumen local & 2/3/6\par Residual: 3 & Cuotas de logs y objetos; monitorear capacidad; límites de retención sin romper replay; backups y restauración aislada; fijar volúmenes a VM-2.\par\textbf{Señal:} Capacidad por encima del margen acordado o error de persistencia.\par\textbf{Rol:} Almacenamiento. \textbf{Verificar:} Medir disco con carga; probar rechazo sin falsa aceptación y recuperar volumen con datos previos. \\[4pt]
\bottomrule
\end{tabularx}
\end{table}


\subsubsection{Evidencias, puertas de avance y cierre de Fase 5}

Si un criterio de hecho falla, no se habilita el paso siguiente. Se registra
la causa, se conserva el estado y se aplica la alternativa nativa. Una falla
de autorización, corrupción, restauración o presión de recursos obliga a
priorizar el monolito antes de reabrir servicios. Los cambios de topología,
contrato o persistencia actualizan conjuntamente ADR, A3, A4 y este plan.

La consigna no asigna F5-E8 ni una prueba P adicional a A5. Se conservan
plan, runbook, matriz CSV y tabla generada, junto a una lista de evidencias
para los ensayos futuros. No se inventa una captura ni una medición.
La generación verifica escalas, productos, IDs y campos de los trece riesgos;
no ejecuta operaciones sobre el servidor.

Con A1--A5 queda documentado el diseño de Fase 5: análisis, decisiones,
contratos, diagramas, orden de migración, reversión y riesgos. Quedan
pendientes la implementación y las pruebas reales de Fases 6 y 7, incluida
la comprobación de cada criterio global y del límite de reversión.
